\section{Framework}

The proposed framework separates forecasting from operational planning decisions. Candidate forecasting models generate demand forecasts. A decision layer then selects or combines these forecasts while accounting for stability, model switching cost, and execution capacity. The selected forecast is converted into a planning signal, such as an inventory target, and evaluated in an operational planning simulator.

The framework is designed to answer whether a forecast that is numerically more accurate also improves executable planning quality. Future experiments will instantiate this framework on public demand datasets and compare direct forecast pass-through policies against stability-aware decision strategies.

The improved decision layer includes feasibility-aware smoothing. Instead of immediately replacing the previous executable plan with a new forecast-driven plan, the planner can adapt gradually with an interpretable smoothing parameter. This mechanism is intended to model execution environments where infrastructure, supplier coordination, staffing, or store operations cannot absorb abrupt period-to-period plan changes.

The framework also includes feasibility-aware ensembles. These strategies combine candidate forecasts with nonnegative weights that sum to one, where the weights can be based on validation forecast accuracy, validation operational loss, or a transparent constrained validation grid. These ensemble variants are included to test whether a soft composition strategy can preserve useful accuracy while reducing execution burden.

Figure~\ref{fig:feasibility-analysis-flowchart} summarizes the planned analysis workflow. The figure is a structural placeholder and does not report empirical findings.

\begin{figure}[htbp]
\centering
\begin{tikzpicture}[
    node distance=0.95cm and 1.05cm,
    stage/.style={
        draw,
        rounded corners=2pt,
        align=center,
        minimum width=2.7cm,
        minimum height=0.82cm,
        fill=blue!5,
        line width=0.45pt
    },
    decision/.style={
        draw,
        diamond,
        aspect=2.0,
        align=center,
        inner sep=1.2pt,
        fill=orange!10,
        line width=0.45pt
    },
    output/.style={
        draw,
        rounded corners=2pt,
        align=center,
        minimum width=2.7cm,
        minimum height=0.82cm,
        fill=green!7,
        line width=0.45pt
    },
    arrow/.style={-{Stealth[length=2.0mm]}, line width=0.45pt}
]
\node[stage] (data) {Operational\\demand data};
\node[stage, right=of data] (forecasts) {Candidate\\forecasts};
\node[stage, right=of forecasts] (signals) {Planning signal\\conversion};
\node[decision, below=of forecasts] (decision) {Feasibility-aware\\decision layer};
\node[stage, left=of decision] (capacity) {Execution capacity\\scenarios};
\node[output, right=of decision] (outputs) {Tradeoff tables\\and figures};

\draw[arrow] (data) -- (forecasts);
\draw[arrow] (forecasts) -- (signals);
\draw[arrow] (signals) -- (decision);
\draw[arrow] (capacity) -- (decision);
\draw[arrow] (decision) -- (outputs);
\end{tikzpicture}
\caption{Feasibility-aware demand-planning analysis workflow.}
\label{fig:feasibility-analysis-flowchart}
\end{figure}


This section is a placeholder. The final version will describe the implemented pipeline, decision layer variants, feasibility-aware smoothing, feasibility-aware ensembles, execution capacity model, and policy drift scenarios in detail.
